\section{The remaining documentary question and source qualifications}

This part records the precise remaining uncertainty and the limits of
the inspected sources. The response to the entire research prompt is
still marked \textbf{INCOMPLETE}: the date of the particular cote~1
trace and genus-zero sheets has not been established. The checked
mathematical theorems retain their explicitly stated hypotheses.

\subsection{The date of the particular Hilbert-space argument}

We have not proved that the crossed-out trace argument on physical PDF
pages~79--80 and the genus-zero lemma on pages~111--112 were written
before 1959. The Montpellier catalogue describes the mathematical
notes as undated and separately lists a letter dated 1953. The
original institutional inventory, reproduced in the 2019 CIRM
presentation, uses the same distinction; it gives no recipient or
individual-sheet identification for that letter
\cite{geoMontpellierCatalogue,algMontpellierInventoryPresentation2019}.

The complete public PDF has 150 physical pages: one modern cover and
149 numbered public pages. Every page has now been displayed at
readable size in the coordinated dating inspection: physical~1--24
by the algebra agent, 25--72 by the geometry agent, 73--110 by the
category agent, and 111--150 by root. These are documentary readings,
not diplomatic transcriptions of every faint or cancelled mark.
No legible date or dated correspondence reference was identified
that ties the trace sheets to a year. The two typed German opening
lines on physical page~26 have no identified date and have not been
identified with the catalogue's 1953 letter.

The three number systems must remain distinct. Physical pages~2, 3
and~4 have consecutive public PDF footers~1, 2 and~3, but archival
pencil numbers~1, 2 and~5. The official description identifies the
lower-left HB-pencil numbering as the archivists' work
\cite[``La num\'erisation de conservation'']{algMontpellierArchiveMethod}.
For the trace argument the correspondence is physical~79--80,
public footers~78--79, archival pencil~80--81; for the lemma it is
physical~111--112, public footers~110--111, archival pencil~112--113.
The skipped pencil numbers~3--4 do not identify their contents.
The archive's general restriction on disseminating correspondents'
letters is a source-access qualification, not proof that those two
numbers are the 1953 letter. We do not infer a date from adjacency,
a numbering gap, or the period when a subject occupied the author.

The dated Dixmier correspondence, institutional inventories for other
early correspondence, and the original 1959 Lidskii endnote have
been checked within the scopes recorded in the source notes. The
24 January 1955 Dixmier letter concerns a determinant inequality;
it does not identify these genus-zero sheets. Lidskii points to the
1955 memoir, and the original Bourbaki expos\'e~91 expressly states
the Hilbert product formula in a text ending with June 1957. These
are dated witnesses to mathematical statements, not dates for the
particular manuscript. The smallest missing witness is a dated
letter, enclosure description, authorial note or archival collation
that identifies these leaves or their distinctive proof with enough
specificity to give a pre-1959 upper bound. The cautious Lead that
the sheets \emph{may} predate 1959 is not replaced by a stronger
priority assertion.

The separately catalogued 1953 thesis is described as 130 folios, but its
original text has not been opened or identified with this dossier
\cite{rootGrothendieckThesisCatalogue1953}. Schwartz's contemporary seminar
introduction explains its dependence on that thesis; the nuclear-operator
lecture's eigenvalue criterion expressly assumes a Hermitian operator
\cite[p.~III]{rootSchwartzSeminarIntroduction1954}
\cite[expos\'e~12, p.~3]{rootSchwartzNuclear1954}. These passages do not
identify the crossed-out general trace proof. The following lecture's
Hilbert-space section uses polar decomposition for the nuclear norm and
again states its eigenvalue criterion for Hermitian operators
\cite[expos\'e~13, \S3, pp.~3--5]{rootSchwartzNuclearCases1954}.
Paumier's precise locator led to the original Bourbaki draft~179,
NBR082 \cite[p.~64, note~73]{rootPaumierSchwartz2014}.
We read its catalogue cover and all eight typed pages. It develops
elementary projective tensor-product properties, with no spectral trace
or genus-zero proof. Its first typed page bears the pencil wording
``Schwartz'' and ``mai53''; the annotation's author and time of inscription
have not been independently authenticated. These marks do not identify
or date the Montpellier sheets \cite{catBourbaki179}. The bounded correspondence,
thesis, course and Brazilian inventory searches are recorded in the
source notes; inaccessible originals remain expressly uninspected.

\subsection{The publication question and edition qualifications}

Proposition~\ref{geo-H-history} records the completed source comparison.
The inspected memoir states both Hilbert genus zero and actual
unordered $\ell^1$ spectral continuity, either sufficient for the
trace formula. Their Hilbert proofs are not supplied in the checked
dependency chains. The complete 1956 Fredholm paper supplies
determinant foundations from which the full trace formula can be
reconstructed; Proposition~\ref{geo-H-1956-reconstruction} gives that
deduction and identifies its added steps. Thus absence of the
manuscript's three-sector argument is not mistaken for mathematical
insufficiency of the published results. The metric \emph{R\'esum\'e}
and Lidskii's original three-page article have also been inspected
and incorporated.

The memoir scan is its fourth printing, dated 1966, with copyright
1955. Line-for-line identity with the first physical printing has
not been established by collation. This is an edition qualification,
not an additional mathematical task invented from the prompt.
The analytic proofs in the checked part do not depend on resolving
the manuscript date.

\subsection{What ``origin of SGA 6'' can mean}

The PS on 23/053--054, the final EGA IV 19.1.8, and SGA 6 VII
1.1--1.7 have been compared with original scans. The mathematical
correspondence is exact, and EGA's printed note refers to the 1966
seminar. The PS itself reports seminar work already under way and a
transitivity argument already checked. Consequently it does not
support a claim that the notion and proofs were first invented in
the act of writing that PS. A stronger claim about its causal
influence on the eventual SGA 6 wording would require a dated
document connecting the drafting stages; that fact is not inferred
from mathematical agreement alone.

\subsection{Scope distinctions, without additional conjectures}

The draft's earlier definition of condition P is absent from the retained
pages. We have therefore proved preservation of both explicitly defined
forms of formal projectivity in the noetherian adic Cohen-lifting theorem.
We do not identify an absent historical definition by assertion.
The natural-bundle theorem states its locality and finiteness hypotheses,
as the requested research package requires. Automatic formal locality
from mere \emph{\'etale} naturality is a stronger statement and is not
used. The programme headings in cote 26 are accompanied by precise
quantified formulations; no universal assertion about unspecified
``structures'' or unspecified moduli problems is invented.

All active states, failed source attempts, and next retrieval steps are
recorded in the dated ledger and the section-specific source logs.
